\chapter{Background and Literature Review}

\section{Traditional Electronic Health Record Systems}
Traditionally, Electronic Health Records (EHRs) are stored in centralized database management systems maintained by individual hospitals or cloud service providers. This centralized approach simplifies data retrieval but introduces significant structural weaknesses. The central database constitutes a lucrative target for malicious actors; a single firewall breach can yield sensitive data belonging to millions of individuals. Moreover, patients typically lack agency and visibility over who accesses their records and for what purpose, raising pervasive ethical and legal concerns.

\section{Blockchain Technology in Healthcare}
Blockchain, initially popularized by cryptocurrencies, has found exceptional utility in enterprise domains where trust, auditability, and decentralization are paramount. In the context of healthcare, blockchain serves as an immutable, append-only ledger that securely logs transactions. Every time a medical file is requested, altered, or uploaded, a transaction is recorded. This cryptographic audit trail heavily deters internal malpractice, as malicious accesses cannot be erased silently. 

However, blockchains inherently duplicate data across all consensus nodes. Storing large diagnostic images or voluminous text reports directly on-chain violates privacy norms and causes severe ledger bloat, culminating in unacceptable performance degradation. This necessitates off-chain data storage models where only pointers (e.g., cryptographic hashes and URIs) are archived on the blockchain.

\section{Hyperledger Fabric}
Hyperledger Fabric is a fundamentally different blockchain framework compared to permissionless systems like Ethereum or Bitcoin. It is a "permissioned" blockchain technology spearheaded by the Linux Foundation. In Fabric, network participants are strictly identified through a Membership Service Provider (MSP), adhering to enterprise security requirements.

Fabric employs a robust execute-order-validate architectural design. Smart contracts (referred to as chaincode) are executed to endorse transactions before they are ordered into blocks. This separation of the ordering service from the execution environment enhances system scalability and mitigates non-determinism. Fabric facilitates the use of standard programming languages (such as JavaScript, Go, and Java) to write chaincode, making it highly suitable for corporate infrastructures.

\section{Confidentiality vs. Integrity}
In a permissioned environment, it is critical to distinguish between confidentiality and integrity. The blockchain ledger naturally enforces data integrity—assuring that historical records have not been maliciously modified. However, the ledger itself does not unconditionally guarantee confidentiality since authorized peer nodes replicate the state. Therefore, strong cryptographic measures must be utilized prior to ledger interaction.

\section{Advanced Cryptographic Constructs}

\subsection{Symmetric Authenticated Encryption (AES-GCM)}
Advanced Encryption Standard (AES) in Galois/Counter Mode (GCM) is an authenticated encryption algorithm. Alongside ensuring data confidentiality through robust symmetric encryption, GCM provides data integrity capability via an Authentication Tag. In our framework, AES-GCM guarantees that if an encrypted file stored in the off-chain object store is tampered with, the decryption phase will explicitly fail rather than return corrupted plaintext.

\subsection{Shamir's Secret Sharing Scheme (SSSS)}
Invented by Adi Shamir, this scheme is a form of threshold cryptography that is mathematically proven to be perfectly secure. A secret $S$ is divided into $n$ parts (or shares). The scheme dictates that any subset of $k$ shares (where $k \le n$) can be utilized to effortlessly reconstruct $S$. Crucially, possessing $k-1$ or fewer shares leaves the secret completely undetermined, preserving informational security. By integrating this threshold scheme, our system prevents scenarios where the compromise of a single node exposes the patient's encryption keys.

\subsection{Node Master Keys (NMK)}
To distribute Shamir shares securely over the network without interception, we utilize the concept of Node Master Keys. A specific share intended for node $A$ is wrapped (encrypted) symmetrically via node $A$'s public/shared key material. This guarantees that even if an attacker intercepts the network traffic or scrutinizes the blockchain metadata, they cannot decipher the underlying share without compromising the specific node hardware.

\section{Role of Large Language Models (LLMs)}
In recent literature, Large Language Models have been proposed as powerful assistants in medical data classification. Rather than relying on simplistic keyword-matching algorithms, LLMs can contextually analyze a clinical text to infer urgency or severity. In this project, an LLM assesses unstructured clinical narratives and categorizes the record's priority. This priority level seamlessly translates into our dynamic threshold strategy—assigning stricter threshold requirements for less critical data, and lighter threshold requirements (e.g., requiring fewer nodes for reconstruction) for highly urgent scenarios, thereby ensuring rapid access during emergencies.
